\section{The tail estimate and the gap}
\label{sec:finish}
The finite comparisons are now in place. We first bound all remaining three-body blocks by a summable multiple of $H_Q$, then identify the zero mode at Laughlin flux.

\begin{lemma}[Three-body tail]\label{lem:tail}
For sufficiently large $Q$,
\begin{equation}\label{eq:tailbound}
 \LL_3\left(W_3\sum_{z=16}^Qq_z(Q)R_{3,z}W_3^\dagger\right)
 \ge-\delta_QH,
\end{equation}
where
\begin{equation}\label{eq:delta}
 \delta_Q=\sum_{\substack{17\le z\le Q\\z\text{ odd}}}
       \frac{d_{3,z}}{d_a}|q_z(Q)|(z+1),\qquad
 \delta_Q\longrightarrow\frac{61}{4096}.
\end{equation}
\end{lemma}
\begin{proof}
For $16\le z\le Q$, the bracket in \eqref{eq:q} is positive, since it is at least $2z-2$. Thus only odd $z$ contribute negatively. A normalized highest vector of $R_{3,z}$ is $\sum_{p=0}^zs_pv_p\otimes e_{z-p}$. Its orbit resolves that spin projector with factor $d_{3,z}$. After wedging and lifting, its quadratic form is
\[
 \norm{\sum_{p=0}^zs_pc_{z-p}(g)B_p(g)\psi}^2
 \le\sum_{p=0}^z\norm{B_p(g)\psi}^2.
\]
Averaging each summand proves \eqref{eq:tailbound} with \eqref{eq:delta}.

For $Q\ge4$, the ratios in $(Q)_z/(2Q-2)_z$ are nonincreasing, so
\[
 \frac{(Q)_z}{(2Q-2)_z}\le\left(\frac Q{2Q-2}\right)^z
 \le(2/3)^z,\qquad \frac{d_{3,z}}{d_a}\le2.
\]
The summands, extended by zero for $z>Q$, are therefore dominated by
$6z(z+1)(2/3)^z$. Dominated convergence and \eqref{eq:qlimit} give
\[
 \lim_{Q\to\infty}\delta_Q
 =\frac32\sum_{z=17,19,\ldots}(3z-1)(z+1)2^{-z}
 =\frac{61}{4096}.
\]
For the last identity, write $z=17+2m$. Then
$(3z-1)(z+1)=900+208m+12m^2$, while the sums of $4^{-m}$,
$m4^{-m}$, and $m^24^{-m}$ are $4/3$, $4/9$, and $20/27$,
respectively. Multiplication by $(3/2)2^{-17}$ gives $61/4096$.
\end{proof}

\begin{proof}[Proof of Theorem~\ref{thm:fock}]
Combine Lemma~\ref{lem:square}, \eqref{eq:Adecomp}, \eqref{eq:A3bound}, Proposition~\ref{prop:A4}, and Lemma~\ref{lem:tail}. They give
\[
 H^2\ge d_a\int K(g)\,dg+
       (1-\eta-\delta_Q-1222\eps-o(1))H.
\]
The integral is positive. The certificate gives
\[
 \eta=\frac{93527408868499}{10^{14}},\qquad\eps=\frac3{10^6},
\]
so the coefficient tends to
\begin{equation}\label{eq:finalmargin}
 1-\frac{93527408868499}{10^{14}}-\frac{61}{4096}
   -1222\frac3{10^6}
 =\gamma_*=\frac{4616733319001}{10^{14}}>\frac1{25}.
\end{equation}
The limit in \eqref{eq:finalmargin} is $\gamma_*$. For each fixed $0<\gamma<\gamma_*$, the coefficient therefore exceeds $\gamma$ for all sufficiently large $Q$. Every convergence threshold is independent of particle number, proving the assertion on $\FF_Q$.
\end{proof}

We now recover the known unique zero mode of Haldane's spherical
parent Hamiltonian~\cite[p.~606, Equations~(6)--(7)]{Haldane}.
The elementary divisibility proof also fixes the exact flux at which the
kernel is a line.

\begin{lemma}[Kernel at Laughlin flux]\label{lem:kernel}
For $N\ge2$ and $Q=3(N-1)$, the kernel of $H_{N,Q}$ is the line spanned by $\Psi_{\mathrm L,N}$.
\end{lemma}
\begin{proof}
Realize $U_Q$ as the homogeneous polynomials of degree $Q$ in spinor coordinates $(u,v)$. Write $b_{ij}=u_iv_j-u_jv_i$. The pair-spin $Q-r$ summand has highest polynomial
\[
 b_{ij}^{\,r}u_i^{Q-r}u_j^{Q-r}.
\]
Since $b_{ij}$ is invariant under simultaneous rotations of the pair, every vector in this summand is divisible by $b_{ij}^r$. Exchange of the pair has parity $(-1)^r$, so antisymmetric pair states have odd $r$.

If $H_{N,Q}\Psi=0$, then, in the ambient tensor space,
\[
 0=\ip\Psi{H_{N,Q}\Psi}=\sum_{i<j}\norm{P_{ij}^{(1)}\Psi}^2.
\]
Thus $P_{ij}^{(1)}\Psi=0$ for every pair. Expanding in the other particles' monomials, the remaining pair sectors all have odd $r\ge3$. Thus $b_{ij}^3$ divides $\Psi$ for every $i<j$. Each bracket $b_{ij}=v_ju_i-u_jv_i$ is primitive as a polynomial in $u_i$, since $v_j$ and $u_jv_i$ are coprime, and is linear over the fraction field of the other variables. Gauss's lemma therefore makes it irreducible in the polynomial ring. Brackets for distinct unordered pairs are nonassociate. They are consequently relatively prime. Unique factorization implies
\[
 \prod_{i<j}b_{ij}^3\mid\Psi.
\]
This product has degree $3(N-1)=Q$ in each particle's spinor variables, exhausting the multidegree of $\Psi$; the quotient is constant.

Conversely, $\Psi_{\mathrm L,N}$ is antisymmetric. Its factor $b_{ij}^3$ leaves pair degrees $Q-3$ in each spinor, whose maximum total pair spin is $Q-3$. Multiplication by the invariant bracket cannot introduce spin $Q-1$. Hence every pair projector annihilates $\Psi_{\mathrm L,N}$.
\end{proof}

\begin{proof}[Proof of Corollary~\ref{thm:main}]
Apply Theorem~\ref{thm:fock} with $\gamma=1/25$. On each particle sector,
the spectral theorem gives
\[
 H_{N,Q}\ge\frac1{25}\bigl(I-P_{\ker H_{N,Q}}\bigr)
\]
once $Q$ exceeds the common threshold. Take $Q=3(N-1)$ and apply
Lemma~\ref{lem:kernel}. Choosing $N_0$ so that $3(N_0-1)$ exceeds the
threshold proves the claim.
\end{proof}

The following comparison specializes the recursive spectral relations of
Lemm, Nachtergaele, Warzel, and Young
\cite[arXiv v1, Theorem~2.1(I), Lemma~2.3, and (2.10)]{LemmNachtergaeleWarzelYoung2026Recursive}.

\begin{corollary}[Charge versus neutral gap at fixed flux]\label{cor:charge-neutral}
For \(N\ge2\), fix \(Q_N=3(N-1)\) and write
\(E_0(n;Q_N)=\min\operatorname{spec}H_{n,Q_N}\) for the coefficient-one
interaction \eqref{eq:physicalH}.  Let \(\operatorname{Gap}_N\) be the
first positive eigenvalue of \(H_{N,Q_N}\), its neutral gap.
The particle-number charge gap at this same flux,
\[
 \Delta_N=E_0(N+1;Q_N)+E_0(N-1;Q_N)-2E_0(N;Q_N),
\]
satisfies
\[
 \Delta_N=E_0(N+1;Q_N)
       \ge\frac{N}{N-1}\operatorname{Gap}_N>0.
\]
In particular, \(\Delta_N\ge N/[25(N-1)]\) for \(N\ge N_0\), with
\(N_0\) as in Corollary~\ref{thm:main}.
\end{corollary}

\begin{proof}
All three sectors are nonempty since \(N+1\le\dim U_{Q_N}=3N-2\).
Lemma~\ref{lem:kernel} gives \(E_0(N;Q_N)=0\) and a one-dimensional
ground space.  The \(N\)-particle sector has dimension greater than one,
so \(\operatorname{Gap}_N>0\).
For a nonzero Laughlin vector \(\Psi\), each pair annihilator \(B_p\)
commutes with every \(c_j\), while
\(\sum_j\|c_j\Psi\|^2=N\|\Psi\|^2\).  Thus some \(c_j\Psi\) is a nonzero
\((N-1)\)-particle zero mode, giving \(E_0(N-1;Q_N)=0\).
The divisibility argument in Lemma~\ref{lem:kernel} excludes a zero mode
in the \((N+1)\)-particle sector: the required product of cubed brackets
would have degree \(3N>Q_N\) in each particle variable.

Let \(P_n\) project onto the lowest-energy eigenspace of \(H_{n,Q_N}\),
and define the overlap operator on the \((N+1)\)-particle sector by
\[
 \mathcal G_N=P_{N+1}\left(\sum_{j=0}^{Q_N}c_j^\dagger P_Nc_j\right)P_{N+1}.
\]
Here \(P_{N+1}\) projects onto a positive-energy ground eigenspace.
For the pure two-body fermionic form \eqref{eq:fockH}, the cited
Theorem~2.1(I), with \(m_0=m_1=2\), gives
\[
 E_0(N+1;Q_N)\ge\frac{N+1}{N-1}E_0(N;Q_N)
   +\frac{N+1-\|\mathcal G_N\|}{N-1}\operatorname{Gap}_N.
\]
The cited Lemma~2.3 gives \(\|\mathcal G_N\|\le\operatorname{rank}P_N=1\).
Since \(E_0(N;Q_N)=0\), this yields
\(E_0(N+1;Q_N)\ge N\operatorname{Gap}_N/(N-1)\).
The two zero energies give the asserted second difference, and
Corollary~\ref{thm:main} supplies the final uniform bound.
\end{proof}
